\section{Effective separation from quantitative analytic bounds}
\label{sec:effective}
The compactness argument proves uniform existence. We now compute
conservative constants from the numerical margins already present in
the class. In particular, no separately calibrated fingerprint order or
separation constant is needed. We also replace derivatives in the
fingerprint by finitely many boundary values. This is a recognition
result for coexisting bridges, not finite-data determination of an
arbitrary analytic body.

\subsection{An explicit continuation estimate}
The following elementary estimate is deliberately nonoptimal. Its role
is to remove an unspecified compactness constant, not to improve the
classical conditioning theory of analytic continuation; compare
\cite{Trefethen}.

\begin{lemma}[A numerical strip modulus]\label{lem:explicit-strip}
Let $f$ be $2\pi$-periodic, holomorphic on $|\operatorname{Im}z|<\rho$,
and bounded there by $B$. Suppose $|f(x)|\le e\le B$ for
$|x|\le a$, where $0<a\le\min(\rho/8,1/2)$. Put
\begin{equation}\label{eq:explicit-strip}
 J=\lceil4/a\rceil,\qquad \vartheta_*=2^{-(J+3)}.
\end{equation}
Then
\[
             \sup_{x\in\R}|f(x)|\le3B(e/B)^{\vartheta_*}.
\]
All constants are uniform under real translation of the observed interval.
\end{lemma}
\begin{proof}
First suppose $0<e<B$ and write $t=e/B$. Interpolate at the $m+1$
equispaced nodes $x_i=-a+2ai/m$, $0\le i\le m$, with $m\ge1$.
For $|z|\le a$ the sum of the absolute Lagrange coefficients is at most
\[
 \sum_{i=0}^m\frac{(2a)^m}{(2a/m)^m i!(m-i)!}
          =\frac{2^m m^m}{m!}\le(2\exp(1))^m<6^m.
\]
Use the circle $|\zeta|=4a$ for the holomorphic interpolation remainder.
Its distance to every node is at least $3a$, and
$|\zeta-z|\ge3a$. Cauchy's formula therefore bounds the remainder by
$(4/3)B(2/3)^{m+1}$. Consequently
\begin{equation}\label{eq:explicit-interpolation}
 \sup_{|z|\le a}|f(z)|
       \le B\{6^mt+(4/3)(2/3)^{m+1}\}.
\end{equation}
Take $m=\lfloor\log(1/t)/\log9\rfloor$ when this integer is positive.
With $\alpha=\log(3/2)/\log9$, the first term is at most $t^\alpha$
and the second at most $(4/3)t^\alpha$. Since $\alpha>1/8$, this
is bounded by $3t^{1/8}$. If $m=0$, use the global bound $B$ and
$t>1/9$ to obtain the same conclusion.

Suppose a disk of radius $a$, centered on the real axis, has bound
$Bu$. Hadamard's three-circles inequality at radii $a,2a,4a$ gives
bound $Bu^{1/2}$ on the middle disk; throughout one may replace
$u$ by $\min(u,1)$. A disk of radius $a$ whose center moves by at most
$a$ is contained in that middle disk. Starting with the first disk,
move in both real directions in steps of at most $a$. At most $J$
steps cover $[-\pi,\pi]$, since $\pi<4$. Each radius-$4a$ disk stays
inside the given strip. The resulting bound is at most
$B(3t^{1/8})^{2^{-J}}\le3Bt^{\vartheta_*}$.
Periodicity finishes the proof. For $e=0$ use the identity theorem,
or take a decreasing positive upper bound for $e$; $e=B$ is immediate.
\end{proof}

\subsection{Numerical geometric constants}
Use fixed physical units. Denote the lower curvature bound by
$\kappa_*>0$. Choose supplied numerical graph bounds $R,M_g>0$ so that
both canonical facing graphs of every bounded bridge extend to $|z|<R$
and satisfy $|\psi_b(z)|\le M_g$ there. Reduce $R$ once, if needed,
so its real part lies inside the permitted local chart neighborhood.
These are quantitative versions of the graph priors in
Section~\ref{sec:setting}. Set
\begin{align}
 r&=R/4,&H&=1+2M_g/R,&u&=\kappa_*r/4,\label{eq:eff-local}\\
 a&=\min\{\rho/8,u/(1+u),1/2\},&
 B&=2\{M+(R_++D_0)3^{\lceil\rho\rceil}\}.\label{eq:eff-strip}
\end{align}
Use $J,\vartheta_*$ from \eqref{eq:explicit-strip}. Define
\begin{align}
 C_*&=3+9(R_++2D_0)/\eta,\label{eq:eff-rigid}\\
 \Delta_*&=\min\{B,\sigma/36,\eta/36,d_0/(4C_*)\}.
                                                        \label{eq:eff-delta}
\end{align}
Choose a positive integer $m_0$ with
$3B2^{-m_0}\le\Delta_*$, and put
\begin{equation}\label{eq:eff-threshold}
 P=m_0 2^{J+3},\qquad E_*=B2^{-P},\qquad \chi_*=E_*/8.
\end{equation}
All these choices use only the stated quantitative priors. Rational
upper and lower enclosures with slack may replace their exact values.
Large exponents should be retained in logarithmic or symbolic form;
rounding $E_*$ to machine zero is not a valid calibration.

\begin{lemma}[From a contact interval to a complete pair]
\label{lem:eff-pair}
Consider two bounded oriented bridges, in their respective canonical
frames. If their gap difference and graph differences on
$[-r/2,r/2]$ give support-value errors at most $E\le B$ on the
normal intervals below, then both complete pair images have Hausdorff
distance at most $3B(E/B)^{\vartheta_*}$ in these frames.
In particular, a gap error $d$ and graph errors at most $v$ give this
conclusion with $E=d+v$.
\end{lemma}
\begin{proof}
Cauchy's estimate gives $|\psi_b'|\le2M_g/R<H$ on $[-r,r]$.
The convex graph has $\psi_b''\ge\kappa_*$ and $\psi_b'(0)=0$.
For $|\theta|\le\arctan(\kappa_*r/4)$, its supporting point is
therefore in $(-r/2,r/2)$. At the source, the support value is
\[
 \max_{|y|\le r/2}\{-\psi_0(y)\cos\theta+y\sin\theta\}.
\]
At the target the corresponding outward normals are
$(-\cos\theta,\sin\theta)$, and its support value is the same
maximum with $\psi_0$ replaced by $g+\psi_1$. Taking maxima is
1-Lipschitz in the uniform norm, so the source support error is at
most $v$ and the target error at most $d+v$. No derivative estimate
is inferred from a small graph-value error. The inequality
$u/(1+u)\le\arctan u$ follows by differentiation for $u\ge0$;
thus the real interval $[-a,a]$ is available at each end.

The Steiner point lies in its convex body. Its distance from a contact
is at most $D_0$; in the common source-contact frame the target center
has norm at most $R_++D_0$. Translating a centered support function
adds a first harmonic whose modulus on the strip is at most the
translation length times $\exp(\rho)$. Since
$\exp(\rho)\le3^{\lceil\rho\rceil}$, the difference of either pair
of support functions is bounded by $B$ on the strip. Rotation or
reflection of the argument preserves that bound. Apply
Lemma~\ref{lem:explicit-strip} at the two normal intervals. For compact
convex bodies the uniform support distance equals the Hausdorff distance.
This proves the assertion, including their relative placement.
\end{proof}

\begin{lemma}[Quantitative rigidity within a single table]
\label{lem:eff-coexisting}
Two oriented bridges of the same table whose canonical source and
target bodies both have Hausdorff distance at most $\Delta_*$ are
translates of one another with the same transverse orientation.
\end{lemma}
\begin{proof}
A support error $\Delta$ changes the Steiner point by at most
$2\Delta$, by its integral formula
$s(C)=\pi^{-1}\int_0^{2\pi}p_C(\theta)(\cos\theta,\sin\theta)\dd\theta$.
Centering thus gives support error at most $3\Delta$, and centered
$L^2$ error at most $3\sqrt{2\pi}\Delta<9\Delta$.
If the source types were different, \eqref{eq:shape} would imply
$\sigma\le9\Delta$, contrary to \eqref{eq:eff-delta}.

Translate the two physical source copies to one fixed copy, with its
Steiner point at zero. Let $U,U'$ be their canonical-frame isometries.
The translation part of $U'^{-1}U$ has norm at most $2\Delta$.
Its orthogonal part cannot be a reflection, by \eqref{eq:refl} and
$9\Delta<\eta$. If it is rotation by the representative angle
$|\phi|\le\pi$, then \eqref{eq:rot} gives
$|\phi|\le9\Delta/\eta$. Every point of either target has norm at
most $R_++2D_0$ relative to its source Steiner point. The target
Hausdorff distance after aligning the physical sources by translation
is consequently at most
\[
 \Delta+2\Delta+(R_++2D_0)|\phi|\le C_*\Delta<d_0.
\]
These targets are actual physical bodies in the same table: translating
a source copy to another source copy uses an element of $\Lambda$.
Distinct physical bodies have distance at least $d_0$, and hence
Hausdorff distance at least $d_0$. The targets must therefore be the
same body. The unique closest segment between this fixed source and
target fixes the canonical normal. The preceding exclusion of reflection
fixes its transverse sign. Thus the bridges are equivalent as asserted.
\end{proof}

\subsection{A computable jet order and a value-only fingerprint}
Choose integers $K\ge2$ and $N\ge2$ satisfying
\begin{equation}\label{eq:eff-integers}
 \frac{2M_g(1/8)^{K+1}}{1-1/8}\le E_*/4,
             \qquad N\ge4Hr/E_*.
\end{equation}
The first may be found by a logarithm bound or integer comparisons;
the second is a ceiling. The finite value fingerprint is
\begin{equation}\label{eq:value-signature}
 S_N=\left(g,\{\psi_b(y_j):b=0,1,\ 0\le j\le N\}\right),
       \qquad y_j=-r/2+rj/N.
\end{equation}
Its reversal $\iota_N$ replaces $j$ by $N-j$ in each block, leaving
the gap fixed. Thus it has exactly the same coherent sign action as
the underlying boundary laws, implemented by a permutation rather than
by differentiation.

\begin{theorem}[Effective finite recognition]\label{thm:effective}
For the explicit constants \eqref{eq:eff-local}--\eqref{eq:eff-integers},
any two inequivalent oriented bridges of a single table satisfy
\begin{equation}\label{eq:effective-separation}
 \|F_K-F'_K\|_\infty>\chi_*\quad\hbox{and}\quad
 \|S_N-S'_N\|_\infty>\chi_*.
\end{equation}
Here $F_K$ uses the normalization $r=R/4$. In particular both
fingerprints separate a bridge from its transverse reversal.
Proposition~\ref{prop:registry} applies to either fingerprint with
error less than $\chi_*/16$ and threshold $\chi_*/2$.
No table-dependent contact values or additional separation oracle enter
the calibration. All integers and thresholds are computable from positive
numerical bounds for the indicated priors.
\end{theorem}
\begin{proof}
Suppose first that the jet fingerprints differ by at most $d\le\chi_*$.
For $|y|\le r/2$, the difference of their degree-$K$ Taylor
polynomials, apart from the separately recorded gap, is at most
$d\sum_{j=2}^K2^{-j}\le d/2$. Cauchy's coefficient bounds give the
sum of the two tails at most the first expression in
\eqref{eq:eff-integers}. Their support errors are therefore at most
$3d/2+E_*/4<E_*$, by Lemma~\ref{lem:eff-pair}.

For the value fingerprint, a nearest grid point is within $r/(2N)$.
The difference of the graphs has derivative bounded by $2H$, so its
uniform norm on the interval is at most $d+Hr/N$. Including the gap,
the two support errors are at most
$2d+Hr/N\le E_*/2$. In both cases Lemma~\ref{lem:explicit-strip}
bounds the complete pair errors by
\[
 3B(E_*/B)^{\vartheta_*}=3B2^{-m_0}\le\Delta_*.
\]
Lemma~\ref{lem:eff-coexisting} proves equivalence, giving the
contrapositive of \eqref{eq:effective-separation}. Reversal is an
inequivalent oriented bridge by asymmetry. The registry proof uses only
an isometric involution and the separated error bands, so applies also
to the permutation $\iota_N$.
\end{proof}

\begin{remark}[Effectivity and size]\label{rem:effectivity}
The input is a list of certified positive numerical bounds, not a
qualitative assertion that some compact analytic class exists. Rational
slack permits all choices above by finite arithmetic and integer
comparisons. For example, if $c_1\ge4Hr/B$ is an integer, then
$N=\max(2,c_1 2^P)$ is valid. The representation $(c_1,P)$ is compact,
whereas writing or measuring every coordinate of that fingerprint may
be prohibitively expensive. The very small exponent
$\vartheta_*$ explicitly records this conditioning. The theorem asserts
neither a sharp fingerprint length nor a polynomial dependence on inverse
analytic, symmetry or separation margins. It supplies a finite recipe
where the preceding compactness proof supplied existence only.
\end{remark}
